% ==============================================================
%  DM_Project_Report.tex  —  Air Quality Assessment in Torino
%  Universita degli Studi di Milano-Bicocca  |  AY 2025-2026
% ==============================================================
\documentclass[11pt,a4paper]{article}

% ---- encoding & language
\usepackage[utf8]{inputenc}
\usepackage[T1]{fontenc}
\usepackage[english]{babel}
\usepackage{lmodern}
\usepackage{microtype}

% ---- layout
\usepackage[top=2.4cm,bottom=2.6cm,left=2.8cm,right=2.8cm]{geometry}
\usepackage{setspace}\setstretch{1.18}

% ---- colours
\usepackage[dvipsnames,svgnames,table]{xcolor}
\definecolor{navyblue}{HTML}{0B1F3A}
\definecolor{accentblue}{HTML}{1D4ED8}
\definecolor{softblue}{HTML}{DBEAFE}
\definecolor{accentred}{HTML}{DC2626}
\definecolor{softred}{HTML}{FEE2E2}
\definecolor{accentgreen}{HTML}{059669}
\definecolor{softgreen}{HTML}{D1FAE5}
\definecolor{accentamber}{HTML}{B45309}
\definecolor{softamber}{HTML}{FEF3C7}
\definecolor{mutedgray}{HTML}{64748B}
\definecolor{bordergray}{HTML}{E2E8F0}
\definecolor{lightbg}{HTML}{F5F7FB}

% ---- tables
\usepackage{booktabs}
\usepackage{tabularx}
\usepackage{array}
\usepackage{colortbl}

% ---- lists
\usepackage{enumitem}
\setlist[itemize]{topsep=4pt,itemsep=3pt,leftmargin=16pt}
\setlist[enumerate]{topsep=4pt,itemsep=3pt,leftmargin=20pt}

% ---- graphics & boxes
\usepackage[most]{tcolorbox}
\usepackage{tikz}

% ---- hyperlinks
\usepackage{hyperref}
\hypersetup{colorlinks=true,linkcolor=accentblue,urlcolor=accentblue,
  pdftitle={Air Quality Assessment in Torino --- DM Report}}

% ---- headers & footers
\usepackage{fancyhdr}
\pagestyle{fancy}\fancyhf{}
\renewcommand{\headrulewidth}{0.4pt}
\fancyhead[L]{\small\color{mutedgray}Air Quality Assessment in Torino}
\fancyhead[R]{\small\color{mutedgray}Data Management Project}
\fancyfoot[C]{\small\color{mutedgray}\thepage}

% ---- section styling
\usepackage{titlesec}
\titleformat{\section}{\large\bfseries\color{navyblue}}{}{0em}{}
  [\vspace{2pt}\color{bordergray}\titlerule[0.6pt]]
\titleformat{\subsection}{\normalsize\bfseries\color{accentblue}}{}{0em}{}
\titleformat{\subsubsection}{\normalsize\bfseries\color{mutedgray}}{}{0em}{}
\titlespacing{\section}{0pt}{16pt}{8pt}
\titlespacing{\subsection}{0pt}{12pt}{5pt}

% ---- caption
\usepackage{caption}
\captionsetup{font=small,labelfont=bf,labelsep=period}

% ---- tcolorbox styles
\tcbset{
  rqbox/.style={
    enhanced,breakable,
    colback=softblue,colframe=accentblue,
    leftrule=4pt,toprule=0.4pt,bottomrule=0.4pt,rightrule=0.4pt,
    left=8pt,right=8pt,top=6pt,bottom=6pt,
    fonttitle=\small\bfseries\color{accentblue},
    sharp corners,
  },
  answerbox/.style={
    enhanced,breakable,
    colback=softgreen,colframe=accentgreen,
    leftrule=4pt,toprule=0.4pt,bottomrule=0.4pt,rightrule=0.4pt,
    left=8pt,right=8pt,top=6pt,bottom=6pt,
    fonttitle=\small\bfseries\color{accentgreen},
    sharp corners,
  },
  infobox/.style={
    enhanced,breakable,
    colback=lightbg,colframe=bordergray,
    boxrule=0.5pt,left=8pt,right=8pt,top=6pt,bottom=6pt,
    sharp corners,
  },
}

% ---- helper: inline tag
\newcommand{\tagitem}[1]{%
  \tcbox[on line,boxrule=0.5pt,colback=softblue,colframe=accentblue,
    left=3pt,right=3pt,top=1pt,bottom=1pt,
    fontupper=\tiny\bfseries\color{accentblue}]{#1}%
}

% ==============================================================
\begin{document}
% ==============================================================

% ---- COVER PAGE
\begin{titlepage}
\color{white}
\begin{tikzpicture}[remember picture,overlay]
  \fill[navyblue] (current page.north west)
    rectangle ([yshift=-7.6cm]current page.north east);
  \node[anchor=north west,xshift=2.8cm,yshift=-1.4cm,
        text=white,font=\tiny\bfseries,text width=14cm]
    at (current page.north west)
    {UNIVERSIT\`A DEGLI STUDI DI MILANO-BICOCCA\enspace\textbf{\textperiodcentered}\enspace DATA MANAGEMENT PROJECT};
  \node[anchor=north west,xshift=2.8cm,yshift=-2.3cm,
        text=white,font=\fontsize{24}{28}\selectfont\bfseries,text width=14cm]
    at (current page.north west)
    {Air Quality Assessment in Torino};
  \node[anchor=north west,xshift=2.8cm,yshift=-4.1cm,
        text=softblue,font=\normalsize,text width=13cm]
    at (current page.north west)
    {Annual PM10 and NO2 analysis across Traffic and Background monitoring stations in
     Torino, enriched with yearly weather indicators and reverse-geocoded station metadata.};
\end{tikzpicture}
\vspace*{7.8cm}
\color{black}
\begin{center}
  \tagitem{EEA DiscoData}\;\tagitem{Open-Meteo API}\;\tagitem{OSM / Nominatim}\;
  \tagitem{BigQuery}\;\tagitem{dlt}\;\tagitem{dbt}\;\tagitem{30 dbt tests passed}
\end{center}
\vspace{1.2cm}
\begin{center}
  {\large\bfseries Data Management Project --- Report}\\[0.5em]
  {\normalsize Course: Data Management and Data Visualization}\\[0.3em]
  {\normalsize Universit\`a degli Studi di Milano-Bicocca (UNIMIB) \textbf{\textperiodcentered} Academic Year 2025--2026}\\[1.2em]
  {\normalsize [Student~1] \quad [Student~2] \quad [Student~3]}\\[0.4em]
  {\small Repository: \href{https://github.com/[repo]}{github.com/[repo]}}\\[0.3em]
  {\small [insert final submission date]}
\end{center}
\end{titlepage}

% ---- ABSTRACT
\begin{tcolorbox}[infobox]
\textbf{Abstract.}~%
This project analyses air quality in Torino by integrating heterogeneous environmental
datasets into a reproducible analytical pipeline. The system combines annual air-quality
statistics from EEA DiscoData, hourly meteorological observations from Open-Meteo, and
reverse-geocoded station metadata from Nominatim/OpenStreetMap. Data are acquired through
automated Python dlt pipelines, stored in Google BigQuery, transformed with dbt, validated
through 30 automated tests, and exposed through SQL queries and a Streamlit dashboard.
The project covers all phases required by the Data Management course: acquisition, storage,
profiling, integration, analysis, and quality improvement.
\end{tcolorbox}

\vspace{0.4cm}
\tableofcontents
\newpage

\begin{abstract}
This project analyzes air quality in Torino by integrating heterogeneous environmental datasets into a reproducible analytical pipeline. The system combines annual air-quality statistics from EEA DiscoData, hourly meteorological observations from Open-Meteo, and reverse-geocoded station metadata from Nominatim / OpenStreetMap. Data are acquired through automated Python pipelines, stored in Google BigQuery, transformed with dbt, validated through automated tests, and exposed through SQL queries and a Streamlit dashboard. The project addresses all phases required by the Data Management course: acquisition, storage, profiling, integration, analysis, and quality improvement. Particular attention is paid to enrichment quality measurement, especially for spatial integration, in line with the course FAQ.
\end{abstract}

\tableofcontents
\newpage

% ==============================================================
\section{Introduction}
% ==============================================================

Air pollution is one of the most significant urban environmental challenges, affecting
public health and urban quality of life. Torino is a relevant study case because the city
hosts both \textbf{Traffic} and \textbf{Background} monitoring stations, enabling direct
comparison of pollutant levels by station type across multiple years.

\subsection{Project objectives}

\begin{enumerate}
  \item Build a reproducible data-management pipeline that acquires, stores, transforms,
        validates, and integrates environmental datasets from multiple sources.
  \item Answer four research-driven questions on PM10 and NO2 patterns through SQL analysis
        and an interactive Streamlit dashboard.
\end{enumerate}

\subsection{Research Questions}

\begin{tcolorbox}[rqbox,title=RQ1 --- Descriptive]
How do PM10 and NO2 annual concentrations differ between Traffic and Background monitoring
stations in Torino during 2022--2024?
\end{tcolorbox}
\begin{tcolorbox}[rqbox,title=RQ2 --- Temporal]
How do annual PM10 and NO2 indicators evolve between 2022 and 2024 for different station types?
\end{tcolorbox}
\begin{tcolorbox}[rqbox,title=RQ3 --- Exceedance analysis]
How do PM10 exceedance days above 50~\textmu g/m\textsuperscript{3} vary across years and station types?
\end{tcolorbox}
\begin{tcolorbox}[rqbox,title=RQ4 --- Integration and enrichment quality]
Can yearly weather indicators and reverse-geocoded station information be integrated into the
analytical warehouse with measurable quality?
\end{tcolorbox}

\subsection{Scope}

The project covers \textbf{4 Torino monitoring stations} for the years \textbf{2022, 2023, and 2024}.
The analytical scope is intentionally descriptive and integrative rather than causal.

% ==============================================================
\section{Data Sources and Acquisition}
% ==============================================================

\begin{table}[h!]
\centering
\caption{Data sources summary}
\small
\begin{tabularx}{\textwidth}{@{}llllr@{}}
\toprule
\textbf{Source} & \textbf{Provider} & \textbf{Method} & \textbf{Raw table} & \textbf{Records} \\
\midrule
Air-quality statistics & EEA DiscoData         & HTTP SQL API / dlt & raw.air\_quality\_statistics\_torino & 669 \\
Hourly weather         & Open-Meteo Historical & REST API / dlt     & raw.weather\_torino\_hourly          & 26\,328 \\
Station geocoding      & Nominatim / OSM       & REST API / dlt     & raw.station\_reverse\_geocoding      & 4 \\
\bottomrule
\end{tabularx}
\end{table}

\subsection{EEA DiscoData}

The core pollutant dataset is extracted from EEA DiscoData AirQualityStatistics via an HTTP SQL
endpoint. The pipeline selects Torino stations with a \texttt{P1Y} (annual) aggregation process
and loads \textbf{669 records} into BigQuery.

\subsection{Open-Meteo Historical API}

Hourly meteorological observations are downloaded from the Open-Meteo Historical Archive
(\textbf{26\,328 records}). These are later aggregated yearly inside the dbt mart
\texttt{fact\_weather\_yearly} to support weather enrichment of air-quality facts.

\subsection{Nominatim / OpenStreetMap}

Reverse geocoding enriches the four monitoring stations with human-readable address fields and
measures the Haversine distance between original coordinates and the geocoded point, enabling
explicit quality assessment of spatial integration.

\subsection{Acquisition constraints}

\begin{itemize}
  \item The three sources use different schemas requiring normalisation at ingestion time.
  \item Nominatim enforces a rate limit and requires an explicit user-agent header.
  \item Spatial enrichment must be evaluated quantitatively, not only executed.
\end{itemize}

% ==============================================================
\section{Storage Layer and Database Design}
% ==============================================================

\subsection{Technology choices}

\begin{tcolorbox}[infobox]
\renewcommand{\arraystretch}{1.25}
\begin{tabularx}{\textwidth}{@{}lX@{}}
  \textbf{Ingestion}      & Python \textbf{dlt} --- declarative pipelines with built-in BigQuery destination \\
  \textbf{Storage}        & \textbf{Google BigQuery} --- serverless, scalable analytical DBMS \\
  \textbf{Transformation} & \textbf{dbt Core} --- versioned SQL models, staging/mart design, declarative tests \\
  \textbf{Dashboard}      & \textbf{Streamlit + Plotly} --- interactive Python-based presentation layer \\
\end{tabularx}
\end{tcolorbox}

\subsection{Warehouse schema}

\begin{table}[h!]
\centering
\caption{Warehouse layers and assets}
\small
\begin{tabularx}{\textwidth}{@{}llX@{}}
\toprule
\textbf{Layer} & \textbf{Asset} & \textbf{Description} \\
\midrule
\rowcolor{lightbg}
Raw     & raw.air\_quality\_statistics\_torino & EEA annual statistics as loaded by dlt \\
Raw     & raw.weather\_torino\_hourly          & Open-Meteo hourly observations \\
\rowcolor{lightbg}
Raw     & raw.station\_reverse\_geocoding      & Nominatim geocoding results \\
Staging & stg\_air\_quality\_statistics\_torino & Cleaned and typed EEA records \\
\rowcolor{lightbg}
Staging & stg\_weather\_torino\_hourly         & Cleaned hourly weather data \\
Staging & stg\_station\_reverse\_geocoding     & Cleaned geocoding records \\
\midrule
\rowcolor{lightbg}
Mart    & dim\_station                         & Station dimension with coordinates \\
Mart    & fact\_air\_quality\_annual           & Annual PM10/NO2 KPIs per station \\
\rowcolor{lightbg}
Mart    & mart\_station\_type\_year\_summary   & KPIs aggregated by station type and year \\
Mart    & mart\_air\_quality\_weather\_yearly  & Air-quality facts enriched with weather \\
\rowcolor{lightbg}
Mart    & dim\_station\_enriched               & Stations enriched with geocoded address \\
Mart    & dq\_reverse\_geocoding\_quality      & Spatial integration quality metrics \\
\bottomrule
\end{tabularx}
\end{table}

% ==============================================================
\section{Data Profiling and Data Quality}
% ==============================================================

\subsection{Automated quality checks with dbt}

\begin{table}[h!]
\centering
\caption{dbt test coverage summary}
\small
\rowcolors{2}{lightbg}{white}
\begin{tabularx}{\textwidth}{@{}lllX@{}}
\toprule
\textbf{Dimension}   & \textbf{dbt test}          & \textbf{Applied to}             & \textbf{Result} \\
\midrule
Completeness         & \texttt{not\_null}         & All key columns in every mart   & \textcolor{accentgreen}{\bfseries PASS} \\
Uniqueness           & \texttt{unique}            & station\_code, composite keys    & \textcolor{accentgreen}{\bfseries PASS} \\
Validity             & \texttt{accepted\_values} & station\_type, enrichment flags  & \textcolor{accentgreen}{\bfseries PASS} \\
Enrichment success   & \texttt{accepted\_values} & weather\_enrichment\_success     & \textcolor{accentgreen}{\bfseries PASS} \\
\midrule
\multicolumn{3}{l}{\textbf{Total}} & \textbf{30 / 30 tests passed} \\
\bottomrule
\end{tabularx}
\end{table}

\subsection{Enrichment quality metrics (FAQ 7 and FAQ 9)}

\begin{tcolorbox}[answerbox,title=Measured quality evidence]
\renewcommand{\arraystretch}{1.3}
\begin{tabular}{@{}ll@{}}
Weather enrichment success        & \textbf{100\%} of station-year records \\
Geocoding station coverage        & \textbf{100\%} (4~/ 4 stations enriched) \\
Average geocoding distance        & \textbf{27.09~m} \\
Maximum geocoding distance        & \textbf{55.34~m} \\
Stations within 50~m              & \textbf{75\%} (3 / 4) \\
Stations within 100~m             & \textbf{100\%} (4 / 4) \\
\end{tabular}
\end{tcolorbox}

\section{Data Integration and Enrichment}

\subsection{Weather Enrichment}

The first integration task links annual air-quality indicators with yearly weather summaries. The resulting mart, \texttt{mart\_air\_quality\_weather\_yearly}, makes it possible to analyze pollutant behavior in the context of meteorological conditions.

The integration is automated through the warehouse pipeline rather than manual spreadsheet joins. This directly satisfies the course requirement that integration must be automated.

\subsection{Spatial Enrichment}

The second integration task enriches station metadata with reverse-geocoded information from Nominatim / OpenStreetMap. The resulting \texttt{dim\_station\_enriched} table contains human-readable address fields and the measured distance between original coordinates and the geocoded point.

This enrichment is particularly aligned with the FAQ suggestion to pay attention to spatial integration cases.

\subsection{Integration Outcome}

The integration outcome is positive on both dimensions:

\begin{itemize}
  \item weather enrichment succeeded for all analytical records
  \item reverse geocoding enriched all stations
  \item the geocoding error remained within a small and interpretable range
\end{itemize}

Therefore, the project satisfies the requirement of integrating at least two datasets and measuring the related enrichment error.

% ==============================================================
\section{Analysis and Results}
% ==============================================================

\subsection{RQ1 --- Traffic vs Background concentrations}

\begin{tcolorbox}[rqbox,title=Research Question 1]
How do PM10 and NO2 annual concentrations differ between Traffic and Background stations?
\end{tcolorbox}

\begin{table}[h!]
\centering
\caption{Annual mean concentrations by station type and year (\textmu g/m\textsuperscript{3})}
\small
\rowcolors{2}{lightbg}{white}
\begin{tabular}{@{}llrrrr@{}}
\toprule
\textbf{Station type} & \textbf{Year} & \textbf{Avg NO2} & \textbf{Avg PM10} & \textbf{Exceedance days} & \textbf{Stations} \\
\midrule
Background & 2022 & 29.35 & 30.21 &  47.5 & 2 \\
Traffic    & 2022 & 44.69 & 35.10 &  71.5 & 2 \\
Background & 2023 & 29.38 & 25.41 &  28.0 & 2 \\
Traffic    & 2023 & 39.16 & 32.23 &  58.0 & 2 \\
Background & 2024 & 27.77 & 22.47 &  39.0 & 2 \\
Traffic    & 2024 & 38.59 & 28.71 &  44.5 & 2 \\
\bottomrule
\end{tabular}
\end{table}

\begin{tcolorbox}[answerbox,title=Direct answer --- RQ1]
Traffic stations show higher annual PM10 and NO2 concentrations than Background stations in
every observed year. In 2024, NO2 is 38.59~\textmu g/m\textsuperscript{3} for Traffic versus
27.77~\textmu g/m\textsuperscript{3} for Background, a gap of 10.82~\textmu g/m\textsuperscript{3}.
\end{tcolorbox}

\subsection{RQ2 --- Temporal evolution}

\begin{tcolorbox}[rqbox,title=Research Question 2]
How do annual PM10 and NO2 indicators evolve between 2022 and 2024?
\end{tcolorbox}

\begin{table}[h!]
\centering
\caption{Year-over-year changes by station type}
\small
\rowcolors{2}{lightbg}{white}
\begin{tabular}{@{}llrrrrrr@{}}
\toprule
\textbf{Type} & \textbf{Year} & \textbf{PM10} & \textbf{NO2} & \textbf{$\Delta$PM10} & \textbf{$\Delta$NO2} & \textbf{$\Delta$PM10\%} & \textbf{$\Delta$NO2\%} \\
\midrule
Background & 2023 & 25.41 & 29.38 & $-$4.80 & $+$0.03 & $-$15.9\% & $+$0.1\% \\
Background & 2024 & 22.47 & 27.77 & $-$2.94 & $-$1.61 & $-$11.6\% & $-$5.5\% \\
Traffic    & 2023 & 32.23 & 39.16 & $-$2.87 & $-$5.53 & $-$8.2\%  & $-$12.4\% \\
Traffic    & 2024 & 28.71 & 38.59 & $-$3.52 & $-$0.57 & $-$10.9\% & $-$1.5\% \\
\bottomrule
\end{tabular}
\end{table}

\begin{tcolorbox}[answerbox,title=Direct answer --- RQ2]
PM10 decreases for both Traffic and Background station types across the period analysed.
NO2 also declines but remains comparatively high at Traffic stations, close to the EU annual
limit of 40~\textmu g/m\textsuperscript{3}.
\end{tcolorbox}

\subsection{RQ3 --- PM10 exceedance days}

\begin{tcolorbox}[rqbox,title=Research Question 3]
How do PM10 exceedance days above 50~\textmu g/m\textsuperscript{3} vary across years and station types?
\end{tcolorbox}

\begin{table}[h!]
\centering
\caption{PM10 exceedance days per station and year (EU limit: 35 days/year)}
\small
\rowcolors{2}{lightbg}{white}
\begin{tabular}{@{}llrrr@{}}
\toprule
\textbf{Station ID} & \textbf{Type} & \textbf{2022} & \textbf{2023} & \textbf{2024} \\
\midrule
IT0469A & Traffic    & 57 & 53 & 34 \\
IT0470A & Traffic    & 86 & 63 & 55 \\
IT1877A & Background & 58 & 37 & 41 \\
IT2168A & Background & 37 & 19 & 37 \\
\midrule
\textbf{Traffic avg}    & & \textbf{71.5} & \textbf{58.0} & \textbf{44.5} \\
\textbf{Background avg} & & \textbf{47.5} & \textbf{28.0} & \textbf{39.0} \\
\bottomrule
\end{tabular}
\end{table}

\begin{tcolorbox}[answerbox,title=Direct answer --- RQ3]
PM10 exceedance days are higher for Traffic stations in all observed years: 71.5 vs 47.5 in 2022,
58.0 vs 28.0 in 2023, and 44.5 vs 39.0 in 2024. Traffic stations exceed the EU maximum of 35
days every year, while Background stations exceed it in 2022 and 2024 but stay below it in 2023.
\end{tcolorbox}

\subsection{RQ4 --- Integration and enrichment quality}

\begin{tcolorbox}[rqbox,title=Research Question 4]
Can yearly weather and reverse-geocoded station information be integrated with measurable quality?
\end{tcolorbox}

\begin{tcolorbox}[answerbox,title=Direct answer --- RQ4]
Yes. Weather enrichment succeeded for all station-year records (100\%) and reverse-geocoded
station enrichment achieved 100\% coverage. The maximum geocoding distance is 55.34~m.
The dbt quality layer completed successfully with 30 tests passed.
\end{tcolorbox}

% ==============================================================
\section{Compliance Summary}
% ==============================================================

\begin{table}[h!]
\centering
\caption{EU and WHO regulatory compliance}
\small
\rowcolors{2}{lightbg}{white}
\begin{tabular}{@{}llll@{}}
\toprule
\textbf{Metric} & \textbf{Limit} & \textbf{Traffic} & \textbf{Background} \\
\midrule
NO2 annual mean      & EU 40\,\textmu g/m\textsuperscript{3} & Near/above (2022)       & Compliant all years \\
PM10 annual mean     & EU 40\,\textmu g/m\textsuperscript{3} & Compliant all years     & Compliant all years \\
PM10 exceedance days & EU $\leq$35 days/year                 & Exceeds all three years & Exceeds 2022 and 2024 \\
NO2 annual mean      & WHO 10\,\textmu g/m\textsuperscript{3}& Exceeds all years       & Exceeds all years \\
PM10 annual mean     & WHO 15\,\textmu g/m\textsuperscript{3}& Exceeds all years       & Exceeds all years \\
\bottomrule
\end{tabular}
\end{table}

% ==============================================================
\section{Reproducibility and Deployment}
% ==============================================================

\begin{table}[h!]
\centering
\caption{Implemented assets}
\small
\begin{tabularx}{\textwidth}{@{}lX@{}}
\toprule
\textbf{Category} & \textbf{Assets} \\
\midrule
dlt pipelines   & open\_meteo.py, eea\_discodata.py, station\_reverse\_geocoding.py \\
Raw tables      & raw.weather\_torino\_hourly, raw.air\_quality\_statistics\_torino, raw.station\_reverse\_geocoding \\
dbt staging     & stg\_weather\_torino\_hourly, stg\_air\_quality\_statistics\_torino, stg\_station\_reverse\_geocoding \\
dbt marts       & dim\_station, fact\_air\_quality\_annual, mart\_station\_type\_year\_summary, \newline
                  fact\_weather\_yearly, mart\_air\_quality\_weather\_yearly, dim\_station\_enriched, dq\_reverse\_geocoding\_quality \\
SQL queries     & annual\_means\_by\_station.sql, annual\_means\_by\_station\_type.sql, pm10\_exceedances\_by\_station.sql \ldots \\
Dashboard       & Streamlit app at \texttt{dashboard\_streamlit\_v2/} \\
Quality         & 30 dbt tests passed \\
\bottomrule
\end{tabularx}
\end{table}

The project is fully reproducible. All environment variables, credentials, and commands required
to run the complete pipeline are documented in the accompanying \textit{Operational Guide}.

% ==============================================================
\section{Limitations and Future Work}
% ==============================================================

\begin{tcolorbox}[infobox]
\textbf{Limitations:}
\begin{itemize}
  \item Annual air-quality indicators limit analysis to year-level trends.
  \item The station scope is limited to the four Torino stations available in EEA DiscoData.
  \item The analysis is descriptive; causal inference would require additional data.
\end{itemize}
\textbf{Future improvements:}
\begin{itemize}
  \item Incorporate daily ARPA data for higher temporal resolution.
  \item Extend the geographic scope to other Italian cities.
  \item Integrate road traffic or land-use spatial datasets.
  \item Track data-quality metrics historically through a monitoring mart.
\end{itemize}
\end{tcolorbox}

% ==============================================================
\section{Conclusion}
% ==============================================================

This project demonstrates a complete end-to-end data-management workflow: three automated
API-based ingestion pipelines, a structured analytical warehouse in BigQuery, dbt transformations
and 30 schema tests, two forms of automated enrichment with measured quality, six SQL analytical
queries, and a Streamlit dashboard aligned with the four research questions.

\textbf{Main findings:}
\begin{itemize}
  \item Traffic stations consistently exhibit worse air-quality indicators than Background stations.
  \item PM10 annual means decrease over 2022--2024 for both station types.
  \item PM10 exceedance days remain above the EU limit at Traffic stations in every observed year.
  \item Weather and spatial enrichment integrate with full coverage and controlled error.
  \item All stations exceed WHO AQG 2021 thresholds for both PM10 and NO2.
\end{itemize}

\end{document}